\chapter{Historical compatibility and future simulation}
\label{ch:math}

\section{Define what is observed before inverting it}

Let $D_k$ be the disclosures available under information cutoff $k$. Each
numeric observation $d_i$ has a context: entity or group perimeter, concept,
measurement basis, dimensions, currency, period and precision. These contexts
are part of the observation operator, not additional columns that can be
discarded after extraction.

A candidate history $h$ contains an opening economic and accounting state,
events through the last historical date $T$, contract terms, accounting
policies, consolidation relationships and latent quantities. Let $M$ denote
the declared model scope: its event templates, admissible rule versions and
allowed unknowns. It excludes arbitrary events that the chosen model has no
way to represent. The reporting operator $O_i(h)$ produces the quantity to
which observation $i$ refers, including its basis and perimeter.

Reported numbers are often rounded. If a number $v_i$ is rounded to the nearest
unit $q_i$, an illustrative admissible interval is
$[v_i-q_i/2,v_i+q_i/2)$ under the corresponding tie convention. This
half-open endpoint choice is used only when that tie rule is declared;
Section~\ref{sec:source-rounding} gives the conservative closed enclosure when
nearest rounding is assumed but the tie rule is unknown. The actual
rounding convention, scale and XBRL accuracy metadata govern the imported
constraint. A disclosure with unspecified precision does not acquire a
convenient tolerance silently. Write $\mathcal{I}_i$ for the justified
observation interval, and $C_M(h)$ for the joint accounting, contractual and
semantic constraints. Then the historical compatibility set is
\begin{equation}\label{eq:histset}
 \Hist_M(D_k)=\{h\in\mathcal{H}_M:
 C_M(h)\ \land\ O_i(h)\in\mathcal{I}_i\text{ for every applicable }i\}.
\end{equation}
Qualitative assertions enter as reviewed logical conditions, rather than
numeric intervals. Incompatible source contexts are never forced into the
same equation.

Equation~\eqref{eq:histset} distinguishes three outcomes. An empty set means
that the observations and declared model are inconsistent. A singleton means
that this model identifies a history at its chosen resolution. A larger set
preserves unresolved alternatives. An empty set does not prove that the
company's accounts are false: a wrong extraction, missing event type or
incorrect perimeter can cause the inconsistency. A singleton does not prove
that the model contains the real history.

\section{Why published statements cannot identify all transactions}

The simplest derivation uses fixed account mappings and a fixed basis. Let
$z\in\mathbb{R}^n$ contain unknown historical movement amounts. For $m$
disclosed aggregates, let $A\in\mathbb{R}^{m\times n}$ be the aggregation
matrix and $y=Az\in\mathbb{R}^m$. If there is a nonzero vector $v$ with
$Av=0$, then
\[
 A(z+v)=Az+Av=y.
\]
Both movement vectors give the same reports whenever both satisfy the remaining
constraints. Thus the feasible family is the affine set $z+\ker A$ intersected
with those constraints. Accounting requirements can shrink this family, but
do not generally eliminate it. For example, cash collections in a receivables
bridge are constrained by
\[
 AR_T=AR_0+\text{credit sales}-\text{collections}
       -\text{write-offs}+\text{other classified movements}.
\]
If the note does not distinguish acquisitions, currency effects and write-offs,
several collections paths may fit the same closing balance.

There is a second ambiguity beyond aggregation. In
Equation~\eqref{eq:running}, identical historical movements arise from different
contractual reset terms. Even recovering $z$ exactly would leave these terms
unidentified. AccDSL therefore infers economically relevant states and
contract terms alongside journal movements. A linear inverse solver over
statement lines alone cannot solve this second problem.

\section{The accounting state and its invariants}

For one entity book with $n_a$ accounts in one functional currency, let
$b_t\in\mathbb{Q}^{n_a}$ be
signed account balances, with debit amounts positive and credit amounts
negative. Exact decimal amounts can be represented as rationals with a power
of ten denominator. A journal batch $j_t$ contains the signed functional-currency
amounts posted by one atomic event or explicitly grouped event batch. Its
balance condition is
\begin{equation}\label{eq:posting}
 \boldsymbol{1}^{\mathsf T}j_t=0,\qquad b_t=b_{t^-}+j_t.
\end{equation}
If $\boldsymbol{1}^{\mathsf T}b_0=0$, substitution gives
$\boldsymbol{1}^{\mathsf T}b_t=
\boldsymbol{1}^{\mathsf T}b_{t^-}+\boldsymbol{1}^{\mathsf T}j_t=0$.
Induction establishes conservation through every committed batch. Opening
balances must therefore be supplied as a balanced, explained entry or a
balanced imported state with unresolved composition explicitly identified.

This proof establishes a bookkeeping invariant. It does not establish that an
amount belongs in the chosen account, that a contract exists, or that an
accounting standard permits a treatment. Posting validation additionally
checks account roles, ownership, currency conversions, applicability of rules
and event preconditions.

The balance vector includes temporary revenue and expense accounts during
the period. Period closing transfers their net balance to equity through
balanced entries. A reporting function can include current profit in equity
before closing, or use the closed book afterwards, but must not include it
twice. Income and cash-flow reports use classified movements and exclude
mechanical closing transfers from operating flows. Consequently statements
are functions of balances, journals, contract state and reporting rules:
\begin{equation}\label{eq:report}
 Y_{[u,v)}=R(b_{v^-},J_{[u,v)},X_{v^-},\rho,[u,v)).
\end{equation}
Here $\rho$ specifies the reporting basis and perimeter, and $X$ contains
economic state beyond ledger balances. Cash-flow classification is part of
$R$, not recoverable from the closing cash balance alone.

Each book has its own functional currency. Quantities such as shares, tonnes
and foreign-currency units are typed separately and never added to this
balance equation. A foreign-currency posting stores its original units and
the functional-currency valuation with its rate source. Translation into a
group's presentation currency and consolidation adjustments occur in a
separate book under explicit rules.

\section{Simulate events from economic state}

For a chosen history $h$, the state at $T$ is $X_T(h)$. It contains balances,
open contracts, lots or cohorts, unpaid accruals, commitments and other
quantities needed for future evolution. Let $s$ identify a scenario definition,
$U_t^s$ its supplied economic drivers, $a_t$ management's actions, and
$\varepsilon_t$ future innovations. A behavioural policy $\pi$ chooses
$a_t$ using the permitted information $\mathcal{F}_t$:
\begin{align}
 a_t&=\pi(X_t,U_{\leq t}^s,\mathcal{F}_t;\theta),\label{eq:actions}\\
 e_{t+1}&=\mathcal G(X_t,a_t,U_{t:t+1}^s,\varepsilon_{t+1};\theta),\label{eq:events}\\
 (j_{t+1},m_{t+1})&=\mathcal A(e_{t+1},X_t,r_{t+1}),\label{eq:account}\\
 X_{t+1}&=\mathcal T(X_t,e_{t+1},j_{t+1},m_{t+1}).\label{eq:transition}
\end{align}
The generator $\mathcal G$ emits events, such as a sale, rate reset, payment or default.
The accounting function $\mathcal A$ applies rule package $r$ and returns journal
postings plus measurement annotations $m$. The transition $\mathcal T$ updates economic
and ledger state atomically. A default can therefore change the loan's legal
status, expected recoveries, recognised allowance and cash collections through
related, but distinct, records.

This ordering is a reference abstraction. Within a date, event dependencies
and phases specify the actual order. Expectations used to measure an allowance
or provision must be calculated from the information permitted at that date.
They cannot use the realised future path of the simulation. An inner
expectation model may use quadrature, analytic formulas or nested simulation;
its approximation and seed belong in the measurement annotation.

Some scenarios specify a deterministic macroeconomic path while leaving
idiosyncratic defaults random. Others supply a joint stochastic driver model.
A risk-neutral measure used to value a derivative is distinct from the
real-world measure used to simulate business outcomes. Every stochastic
adapter declares its measure and conditioning information. The accounting
function receives the resulting valuation with that declaration.

\section{Preserve sets and distributions separately}

Let $K_{h,s}$ be the probability law of future reports produced by the forward
model conditional on history $h$, scenario $s$ and its declared behavioural
assumptions. A deterministic simulation is the special case of a point mass.
When no probabilities over historical explanations are justified, the natural
answer is a family of laws,
\begin{equation}\label{eq:family}
 \mathcal{K}_{M,s}(D_k)=\{K_{h,s}:h\in\Hist_M(D_k)\}.
\end{equation}
The system can report each explanation, bounds on a deterministic outcome,
or lower and upper event probabilities across this family. Those are different
result types. A range of conditional means is not a confidence interval, and
an envelope of conditional quantiles is not itself a single predictive
distribution.

If a defensible probability measure $\mu(dh\mid D_k,M)$ supported on
$\Hist_M(D_k)$ is supplied, histories and report outcomes must have specified
measurable spaces. Assume that $K_{h,s}$ is a probability kernel: it is a
probability measure on report outcomes for each $h$, and $h\mapsto K_{h,s}(B)$
is measurable for each measurable event $B$. The predictive law is then
\begin{equation}\label{eq:mixture}
 \Prob(Y^s\in B\mid D_k,M)=
 \int K_{h,s}(B)\,\mu(dh\mid D_k,M).
\end{equation}
For discrete explanations this is the usual weighted sum. For continuous
unknowns, integration uses the declared measure over those unknowns, not a
count of solver solutions. Bayesian weights require a prior and a likelihood;
expert weights remain labelled as elicited. A likelihood must not treat a
statement total and all its components as independent repeated evidence.
Deterministic equality constraints require a reduced parameterisation or
proper conditional-measure treatment, rather than an undocumented Gaussian
penalty with near-zero variance.

Suppose the fixed-loan explanation receives weight $w$ and the floating-loan
explanation receives $1-w$, solely as an illustration. Under the rate shock,
\[
 \Prob(I=5)=w,\qquad\Prob(I=7)=1-w,
 \qquad\E[I]=7-2w.
\]
The expected expense is useful but hides the two different outcomes. For a
general scalar output with finite second moment, put $m(h)=\E[Y\mid h]$.
Write $Y-\E[Y]=(Y-m(h))+(m(h)-\E[Y])$. The cross term vanishes because
$\E[Y-m(h)\mid h]=0$. Taking expectations of the squared expression gives
\begin{equation}\label{eq:variance}
 \operatorname{Var}(Y)=
 \E_h[\operatorname{Var}(Y\mid h)]+
 \operatorname{Var}_h(\E[Y\mid h]).
\end{equation}
The first term is uncertainty within a chosen explanation; the second captures
differences between explanations. This decomposition requires the measure
$\mu$. Without it, the system reports the family in
Equation~\eqref{eq:family}, rather than manufacturing a variance.

Weights themselves can be uncertain. For an admissible class $\mathcal{W}$
of measures, AccDSL can compute
\[
 \underline p(B)=\inf_{\mu\in\mathcal{W}}\int K_{h,s}(B)\,d\mu(h),
 \qquad
 \overline p(B)=\sup_{\mu\in\mathcal{W}}\int K_{h,s}(B)\,d\mu(h).
\]
If every measure supported on the feasible histories is allowed and point
masses are included, these bounds reduce to the infimum and supremum of
$K_{h,s}(B)$ over histories: averages cannot lie outside the range of their
integrands, and point masses attain, or approximate, the extremes. Narrower
weight restrictions require an explicit justification and a separate
optimisation problem.

\section{Resolution and computational limits}

For a requested horizon $H$ and scenario family $\mathcal{S}$, histories may be
grouped only when they have the same conditional law for all requested
outputs:
\[
 h_1\sim_{H,\mathcal{S}}h_2
 \quad\Longleftrightarrow\quad
 K_{h_1,s}=K_{h_2,s}\text{ on those outputs for every }s\in\mathcal{S}.
\]
This definition explains why individual sales can sometimes be replaced by
cohorts, while fixed and floating loans usually cannot be merged for an
interest-rate stress. Exact equivalence can be difficult to prove. A grouping
justified only by approximate numerical agreement carries an error criterion
and is not labelled exact. Source provenance remains available even when
execution shares a common state.

Bounds require comparable honesty. A linear program can provide certified
extrema for a supported linear model. A set of feasible sampled histories
provides observed extrema: its minimum is no smaller than the true infimum,
and its maximum is no larger than the true supremum. It therefore supplies
an inner exploration of possibilities, not a certified outer bound. Interval
methods may provide outer bounds that no single history attains. Every result
records which of these meanings applies.

The initial implementation offers exact bounded choices and linear constraints,
symbolic continuous families such as $\alpha\in[0,1]$, and explicit scenario
execution at selected feasible states. General nonlinear partial identification,
automatic model discovery and exhaustive enumeration of all economic histories
are research extensions, not promises made by the core language.

\section{A solver formulation for ordinary rounded disclosures}

The debt example used an exact principal. Real disclosures often do not.
Suppose debt is shown as 100 and interest as 5, each rounded to the nearest
unit with the declared half-open intervals $D\in[99.5,100.5)$ and
$I_0\in[4.5,5.5)$. Using an unknown fraction $\alpha$ and unknown total $D$
would introduce the product $\alpha D$. That expression is bilinear and is
outside the promised linear fragment.

Instead use principal amounts $F$ and $V$. The compatible family satisfies
\begin{equation}\label{eq:rounded-debt}
 F,V\geq0,\quad 99.5\leq F+V<100.5,\quad
 4.5\leq0.05(F+V)<5.5.
\end{equation}
Future interest is $I_1=0.05F+0.07V$. Its infimum is
$0.05\times99.5=4.975$, attained with $F=99.5,V=0$. Its supremum is
$0.07\times100.5=7.035$, approached as $F=0$ and $V$ tends to 100.5
from below, but not attained. A closed interval is a valid enclosure, but
its upper endpoint must not be reported as attainable. The exact solver check separately
tests that no better value exists and whether each endpoint is attainable.

For $D>0$, the amount representation is equivalent to the fraction
representation through $D=F+V$ and $\alpha=F/D$. At zero principal the fraction
is undefined and must be omitted or treated as irrelevant, not recovered.
A known fraction restriction $\alpha\geq a$ becomes $F\geq aD$ when
$D>0$ and $a$ is fixed. If both $a$ and $D$ are unknown, that product remains
nonlinear. Reformulation has a domain of validity; it is not permission to
send arbitrary formulas to a linear solver.

The supported exact fragment contains rational linear equalities, weak and
strict inequalities, and finite Boolean or enumeration choices. Guard
partitions become finite disjunctions. Posting in integer cents introduces a
different integer problem if performed before bounding. Consequently the
request specifies whether bounds apply before quantisation, after a monotone
declared rounding map, or to a bounded integer model. Approximate numerical
valuations carry error metadata outside the exact rational kernel. Solver
timeout, unsupported theory and unbounded objective have distinct responses;
none is relabelled as infeasibility.

\section{Building the family with explicit templates}

The reconstruction input is a model-scope record, not an invitation to invent
unlimited journals. It lists event templates, slots per entity and period,
allowed contract classes, source mappings, opening-state domains, calendars,
resolution and excluded operations. A template gives an incidence vector for
postings and a transition for economic state. Northstar's annual bridges
arise by substituting its collection, purchase and payment templates into
opening-plus-movement identities.

For each allowed slot, create an activity flag and nonnegative amount; an
inactive slot has zero amount. Instantiate the template's posting and state
constraints, unify variables referring to the same observation or economic
contract, and match report expressions to source domains through validated
context mappings. Add cross-note restrictions and choice-group dependencies
to this same formula. No separate account-by-account reconstruction is allowed.
Mark every constraint with its evidence, rule or assumption origin.

Enumerate the finite structural choices in a deterministic order. For each
cell, ask for feasibility and retain the continuous formula, rather than
enumerating real-valued solutions. An infeasible cell preserves a sufficient
conflicting subset; an undecided cell preserves its formula and reason. Search
is complete within the declared scope only when every finite cell has been
decided, every applicable observation has been matched or explicitly qualified,
and no resource truncation remains. A witness establishes existence in one
cell. It does not establish complete search, and complete search does not
establish that the scope includes the company's real history.

In Northstar the declared templates identify annual collections, purchases
and supplier payments, but retain every $F\in[0,100]$. Adding a verified note
that $F=60$ narrows this family by conjunction. Adding a proposed ``all floating''
reading to that revised view produces a rejected cell with the conflicting
note attached. Adding a rule that expenses a capital contribution would instead
fail treatment admission: inconsistent evidence and unsupported accounting
have different diagnoses.

\section{Which details may safely be combined}

A resolution record names the outputs, horizon, scenario family, grouping key,
required state, transformation and exact or approximate status. Safe exact
transformations are deliberately few. Cash balances in the same book and
currency can be added. Homogeneous inventory lots may be combined under a
weighted-average policy if ownership, cost basis, impairment state and future
costing treatment coincide. Debt principal can be added only within a group
sharing currency, coupon/reset function, fixing and payment calendars,
maturity, options, guarantees and measurement method. The transition and
reporting maps must be additive on that group.

The last condition supplies a short proof obligation. If a transition obeys
$T(x_1+x_2,u)=T(x_1,u)+T(x_2,u)$ and reporting is additive for every admitted
driver $u$, induction preserves equality after grouping. This does not apply
to covenant thresholds, counterparty netting, nonlinear guarantees or
lot-specific impairment. Two fixed loans with the same dates and rate can
share an execution state under this test. The fixed/floating pair cannot;
neither can two policies with different guarantees merely because their
current reserves match.

A changed horizon, scenario, rule version or requested output invalidates the
resolution approval and reruns the sufficiency check. Extending beyond a debt
maturity requires schedule detail that an earlier short-horizon view could
omit. Approximate cohorts retain the detailed source references and a stated
error criterion tied to outputs. Unless that error is bounded or evaluated,
the grouping remains an approximation awaiting evidence, not an exact quotient
of the history family.

\section{A finite probability model that can actually be evaluated}

Optional probability weighting begins with disjoint economic hypotheses and
a declared joint measure. As an illustration, assign fixed and floating
hypotheses prior probabilities 0.6 and 0.4. The stress itself leaves these
weights unchanged and gives mean interest $0.6(5)+0.4(7)=5.8$.
Suppose genuinely new observed evidence has likelihood 0.2 under fixed and
0.8 under floating. The unnormalised posterior masses are 0.12 and 0.32;
normalisation by 0.44 gives $3/11$ and $8/11$. These likelihoods are
illustrative inputs, not extraction-confidence scores.

The finite reference algorithm enumerates joint hypothesis and innovation
states, evaluates their declared joint mass and any observation likelihood
once, checks the nonzero normalising constant, and propagates the entire
statement vector on each state. It sums masses only across disjoint events.
Two overlapping labels such as ``fixed share at least one half'' and ``all
fixed'' are query views and cannot become two independent mixture components.
Totals and their components share an observation block so their evidence is
not counted twice. Continuous inference needs an additional parameterisation,
density or conditional measure and an assessed numerical method; the finite
algorithm supplies no guarantee for those extensions.

Under the prior, before the new observation is admitted, if the floating branch
has an unsupported valuation operation, its mass 0.4
remains unresolved. The run may report the fixed result with successful mass
0.6, but cannot report 5 as the unconditional mean by renormalising away the
failure. A separately requested mean conditional on successful execution must
show that conditioning event and its mass. For unweighted families, an
unresolved branch blocks a complete envelope over the family. Economic default
is retained as a model outcome with remaining obligations, rather than being
treated as a computational error to remove.

\section{Numerical precision is part of a financial method}
\label{sec:numerical-methods}

The affine family is useful for Northstar because the scenario fixes prices,
rates and operating coefficients while the debt share remains unknown. Other
contracts require roots or discount factors. Their numerical treatment must
state what is exact, what is enclosed and when a posted amount is determined.
An accounting name on a library call does not make its result exact.

Consider positive payments $C_i$ at positive times $t_i$ and a positive price
$P$. The effective annual rate solves
\[
 f(r)=\sum_i\frac{C_i}{(1+r)^{t_i}}-P=0,\qquad r>-1.
\]
Each term discounts a payment to the valuation date. In this positive-payment
case,
\[
 f'(r)=-\sum_i\frac{t_i C_i}{(1+r)^{t_i+1}}<0.
\]
Consequently a certified sign-changing bracket contains exactly one root.
The method uses stated bracket bounds and an expansion cap, evaluates signs
with outward error bounds, and bisects until the requested rate or value
enclosure is reached. An uncertified sign, unavailable bracket or exhausted
accuracy budget yields an unresolved numerical result. Mixed-sign cash flows
can have several internal rates of return; the preceding monotonicity argument
does not apply and cannot justify returning the first Newton iterate that
converges.

Net proceeds 100 and a payment 121 in two years give exactly 10\%. An immediate
fee of 2 changes proceeds to 98 and changes that root. Annual-effective 21\%
for half a year has the exact accumulation factor $\sqrt{1.21}=1.1$, so interest
on 100 is 10. Simple proration would give 10.5 and computes a different contract.
At annual-effective 10\%, the half-year factor $\sqrt{1.1}$ is irrational. For
rational time $p/q$ and positive $1+r$, exact rational power comparisons can
enclose the positive $q$th root of $(1+r)^p$. A negative $p$ first requires the
reciprocal. General exponential discounting needs a separately specified
error-controlled method.

A posting quantum of 0.01 does not eliminate this uncertainty. The whole
certified monetary enclosure must lie in one rounding cell under the admitted
tie rule. An enclosure from 1.0049 to 1.0051 crosses a half-cent boundary and
does not establish one posted amount. Refine within a declared budget or retain
rounding unresolved. This rule separates evidence for a numerical answer from
the accounting authority to use that answer.

\section{Updating a law without counting evidence twice}
\label{sec:inference}

Two financial histories can be compatible with the same reports without being
equally probable. Once an explicitly justified law is supplied, genuinely new
evidence may change it. For finite histories $h$, prior mass $p(h)$ and an
admitted observation likelihood $L(D\mid h)$, Bayes' rule gives
\[
 p(h\mid D)=\frac{p(h)L(D\mid h)}{\sum_g p(g)L(D\mid g)}.
\]
The denominator is the probability of that evidence under the declared model.
If it is zero, the condition is impossible in that model; dropping incompatible
histories and renormalising an empty collection cannot repair it.

For two histories with assumed prior masses $1/3$ and $2/3$ and likelihoods
$3/4$ and $1$, the unnormalised masses are $1/4$ and $2/3$. Their sum is $11/12$,
so the posterior masses are $3/11$ and $8/11$. If the second history's cash
calculation fails, the first still has mass $3/11$ and unresolved cash retains
mass $8/11$. A report giving the successful history probability one would
answer a different, success-conditioned question.

An observation has a semantic information identity as well as a file locator.
Reimporting the same observation through another file does not create a second
measurement. Repeating an identical admitted condition is idempotent: it keeps
the law and information lineage unchanged. Submitting an already consumed
observation as a fresh likelihood block is rejected as duplicate information.
Equal selected member sets alone do not prove that two conditions have the same
meaning. An admitted conditional kernel may divide a region's existing mass
when it is refined; the number of new representatives never determines that
division.

Continuous inference adds a measure choice. One must specify coordinates $z$,
domain, base measure, prior, observation mechanism and any transformation
Jacobian. \citet[Section 2.2, equations (2.5)--(2.8)]{stuart2010} supplies the
inverse-problem formulation, but does not establish identification or easy
sampling for a mixed accounting model. In particular, an exact accounting
equality can confine a model to a lower-dimensional set. Multiplying an ordinary
Lebesgue density by its zero-volume indicator is not a posterior construction.
One must parameterise the constrained model with its declared measure or supply
an actual measurement-error model. Adding convenient small noise changes the
target.

\section{A continuous example whose posterior can be checked}

A linear Gaussian model supplies a tractable comparison for later numerical
inference. Assume the unknown economic quantity has a Gaussian prior
$z\sim\mathcal{N}(m_0,V_0)$ with positive definite covariance $V_0$.
Let the observation be $y=Hz+\epsilon$, where
$\epsilon$ is independent Gaussian noise with positive definite covariance
$R$. The measurement-error assumption is substantive: it describes this example,
not rounding in every financial statement.

Multiplying prior and likelihood adds their quadratic forms. Terms that depend
on $z$ are
\begin{align*}
 &(z-m_0)^\top V_0^{-1}(z-m_0)
 +(y-Hz)^\top R^{-1}(y-Hz)\\
 &\quad=z^\top(V_0^{-1}+H^\top R^{-1}H)z
 -2z^\top(V_0^{-1}m_0+H^\top R^{-1}y)+\text{constant}.
\end{align*}
Completing the square therefore gives posterior covariance $V$ and mean $m$
through
\[
 V^{-1}=V_0^{-1}+H^\top R^{-1}H,\qquad
 V^{-1}m=V_0^{-1}m_0+H^\top R^{-1}y.
\]
The prior and observation contribute precision and precision-weighted
information. An implementation should solve factored linear systems rather
than explicitly form these inverse matrices, and reject an invalid covariance
or unsupported rank condition.

For $m_0=0$, $V_0=4$, $H=1$, $R=1$ and $y=2$, posterior precision is
$1/4+1=5/4$, variance is $4/5$ and mean is $(4/5)2=8/5$. A conditional cash
quantity $100+10z$ has mean 116 and variance 80. A second \emph{independent}
measurement of 2 would give precision $9/4$, variance $4/9$ and mean $16/9$.
Reusing the first measurement would produce those same numbers but would be
wrong for the stated information. The analytic example can distinguish numerical
error from information double counting.

For a bounded non-Gaussian scalar model, quadrature requires a stated domain
or tail truncation and sufficient integrand error bounds to enclose the
normalizer and requested moments. Without a positive lower normalizer bound,
the ratio remains unresolved. Mixed models can enumerate finite choices and
integrate each conditional continuous model, but combining branches requires
their evidence factors, not draw counts. Optional MCMC needs target-specific
parameterisation, tuning, independent chains, tractable correctness comparisons,
convergence and Monte Carlo error checks. This proposal supplies no generally
validated sampler or universal tolerance.

Nor can a more elaborate inference engine identify a quantity that the
observation model never sees. If reports constrain only total debt 100, the
likelihood may be constant in its floating share $x$. Its conditional prior
then remains unchanged. A uniform prior is a chosen assumption, not a discovery
from those reports. The inspected working version of
\citet[Sections 2.1--2.3, equations (2.2)--(2.4)]{giacomini2021} motivates
retaining posterior classes where a conditional prior cannot be revised.
For Northstar there is a simple direct deduction: under all probability
measures on $[0,1]$, the expectation of $5+2x$ lies in $[5,7]$, with endpoints
attained by point masses. This deduction imports no asymptotic theorem from
that paper's structural-vector-autoregression setting.

\section{A stress is not an observed condition or an intervention}

The interest-rate stress in Northstar fixes a future driver path. It supplies
no new evidence about the historical debt share. An observational forecast
instead conditions a joint law on information. A causal intervention needs
structural equations and an exogenous disturbance law, because it replaces
the mechanism determining a chosen variable while leaving other mechanisms
in place \citep[Section 3.2, equations (5)--(9) and (15)--(17)]{pearl2009}.

An elementary example makes the difference numerical. Let $U$ equal zero or
one with probability $1/2$ each, and assume
\[
       X=U,\qquad Y=2X+U.
\]
Observing $X=1$ establishes $U=1$ and hence $Y=3$. Intervening to set $X=1$
replaces the first equation. It leaves the original law of $U$, so $Y$ is 2 or
3 with equal mass and mean 2.5. Selection of a high-investment company and an
imposed investment policy can therefore answer different questions even under
one fully specified model. Financial filings alone do not establish these
causal equations or identify their effects.

An intervention request must name its target, replacement equation, structural
parents and disturbance law and satisfy acyclicity or a separately supported
equilibrium rule. Correlated disturbances remain joint. The initial company
profile supports imposed exogenous stress; the continuous and causal methods
described here require separately admitted extension schemas and signatures.
They do not silently expand the initial finite-law interface.

At every decision date, behavior and measurement receive the admitted historical
view and permitted simulated past. The scheduler's knowledge of a complete
future stress path is not permission for a manager to foresee it. A perfect-
foresight policy needs an explicit additional assumption. Common random numbers
can improve a paired numerical comparison while establishing no causal effect
in actual company data.

